% GENERATED FILE. Edit canonical lessons, then run npm run generate:books.
% canonical-source-hash: fnv1a64:ce2d5be79af820ee
% canonical-lessons: TA-C227-tiruttu, TA-C227-vicari, TA-C227-pativu-cey, TA-C227-munpativu-cey, TA-C227-vinnappi

\chapter{To correct, To enquire, To register, To book in advance, To apply}
\label{ch:ta-to-correct-to-enquire-to-register-to-book-in-advance-to-apply}
\hlchaptermodality{227}

\begin{chapteropening}
\textbf{By the end of this chapter:} \emph{I can say to correct, to enquire, to register, to book in advance and to apply.}

Complete the last lesson of chapter 227: I can say to correct, to enquire, to register, to book in advance and to apply.
\end{chapteropening}

\section[\texorpdfstring{\ta{திருத்து} (tiruttu)}{திருத்து (tiruttu)}]{\ta{திருத்து} (tiruttu) — to correct}

\label{lesson:TA-C227-tiruttu}



\begin{quote}
Before the new one: say the Tamil for to swallow, then the Tamil for to chew.
\end{quote}

\subsection*{You'll want to know: \ta{திருத்து}}
\textbf{\ta{திருத்து}} — \emph{tiruttu} — \textquotedblleft{}to correct\textquotedblright{}.

\begin{cousinweb}
\textbf{\ta{திருத்து}} (\emph{tiruttu}) is native Dravidian. DEDR 3251 lists Malayalam \emph{tiruttuka}, Kannada \emph{tiddu} and Telugu \emph{diddu}. It is the causative of \emph{tiruntu}, \textquotedblleft{}to be correct\textquotedblright{}.
\end{cousinweb}

\begin{grammarlens}[title={What you've built}]
One more word for this chapter.
\end{grammarlens}

\subsection*{Your turn}
\begin{itemize}
\raggedright
  \item \emph{Say it:} \emph{tiruttu}
  \item \emph{Say it:} \emph{tiruttu}, once more
  \item \emph{Say it:} say \emph{tiruttu}
  \item \emph{Recall:} say the Tamil for to swallow, then the Tamil for to chew, then say \emph{tiruttu} again
\end{itemize}

\begin{tcolorbox}[breakable,colback=teal!4,colframe=teal!35!black,title={Before you move on}]
What does \ta{திருத்து} mean? (\textquotedblleft{}To correct\textquotedblright{}.) Say it once more.
\end{tcolorbox}
\section[\texorpdfstring{\ta{விசாரி} (vicāri)}{விசாரி (vicāri)}]{\ta{விசாரி} (vicāri) — to enquire, to ask about}

\label{lesson:TA-C227-vicari}



\begin{quote}
Before the new one: say the Tamil for to chew, then the Tamil for to correct.
\end{quote}

\subsection*{You'll want to know: \ta{விசாரி}}
\textbf{\ta{விசாரி}} — \emph{vicāri} — \textquotedblleft{}to enquire, to ask about\textquotedblright{}.

\begin{cousinweb}
\textbf{\ta{விசாரி}} (\emph{vicāri}) is a verb made from Sanskrit \emph{vicāra}, \textquotedblleft{}pondering, enquiry\textquotedblright{} (CDIAL 11683).
\end{cousinweb}

\begin{grammarlens}[title={What you've built}]
One more word for this chapter.
\end{grammarlens}

\subsection*{Your turn}
\begin{itemize}
\raggedright
  \item \emph{Say it:} \emph{vicāri}
  \item \emph{Say it:} \emph{vicāri}, once more
  \item \emph{Say it:} say \emph{vicāri}
  \item \emph{Recall:} say the Tamil for to chew, then the Tamil for to correct, then say \emph{vicāri} again
\end{itemize}

\begin{tcolorbox}[breakable,colback=teal!4,colframe=teal!35!black,title={Before you move on}]
What does \ta{விசாரி} mean? (\textquotedblleft{}To enquire, to ask about\textquotedblright{}.) Say it once more.
\end{tcolorbox}
\section[\texorpdfstring{\ta{பதிவு} \ta{செய்} (pativu cey)}{பதிவு செய் (pativu cey)}]{\ta{பதிவு} \ta{செய்} (pativu cey) — to register, to record}

\label{lesson:TA-C227-pativu-cey}



\begin{quote}
Before the new one: say the Tamil for to correct, then the Tamil for to enquire.
\end{quote}

\subsection*{You'll want to know: \ta{பதிவு} \ta{செய்}}
\textbf{\ta{பதிவு} \ta{செய்}} — \emph{pativu cey} — \textquotedblleft{}to register, to record\textquotedblright{}.

\begin{cousinweb}
\textbf{\ta{பதிவு} \ta{செய்}} (\emph{pativu cey}) is native: \emph{pativu}, \textquotedblleft{}impression, record\textquotedblright{}, from \emph{pati}, \textquotedblleft{}to be imprinted\textquotedblright{} (DEDR 3911), and \emph{sey}, \textquotedblleft{}to do\textquotedblright{} (DEDR 1957): \textquotedblleft{}make a record\textquotedblright{}.
\end{cousinweb}

\begin{grammarlens}[title={What you've built}]
One more word for this chapter.
\end{grammarlens}

\subsection*{Your turn}
\begin{itemize}
\raggedright
  \item \emph{Say it:} \emph{pativu cey}
  \item \emph{Say it:} \emph{pativu cey}, once more
  \item \emph{Say it:} say \emph{pativu cey}
  \item \emph{Recall:} say the Tamil for to correct, then the Tamil for to enquire, then say \emph{pativu cey} again
\end{itemize}

\begin{tcolorbox}[breakable,colback=teal!4,colframe=teal!35!black,title={Before you move on}]
What does \ta{பதிவு} \ta{செய்} mean? (\textquotedblleft{}To register, to record\textquotedblright{}.) Say it once more.
\end{tcolorbox}
\section[\texorpdfstring{\ta{முன்பதிவு} \ta{செய்} (muṉpativu cey)}{முன்பதிவு செய் (muṉpativu cey)}]{\ta{முன்பதிவு} \ta{செய்} (muṉpativu cey) — to book in advance, to reserve}

\label{lesson:TA-C227-munpativu-cey}



\begin{quote}
Before the new one: say the Tamil for to enquire, then the Tamil for to register.
\end{quote}

\subsection*{You'll want to know: \ta{முன்பதிவு} \ta{செய்}}
\textbf{\ta{முன்பதிவு} \ta{செய்}} — \emph{muṉpativu cey} — \textquotedblleft{}to book in advance, to reserve\textquotedblright{}.

\begin{cousinweb}
\textbf{\ta{முன்பதிவு} \ta{செய்}} (\emph{muṉpativu cey}) is native: \emph{muṉ}, \textquotedblleft{}before\textquotedblright{} (DEDR 5020), \emph{pativu}, \textquotedblleft{}record\textquotedblright{} (DEDR 3911), and \emph{sey}, \textquotedblleft{}to do\textquotedblright{} (DEDR 1957): \textquotedblleft{}record in advance\textquotedblright{}.
\end{cousinweb}

\begin{grammarlens}[title={What you've built}]
One more word for this chapter.
\end{grammarlens}

\subsection*{Your turn}
\begin{itemize}
\raggedright
  \item \emph{Say it:} \emph{muṉpativu cey}
  \item \emph{Say it:} \emph{muṉpativu cey}, once more
  \item \emph{Say it:} say \emph{muṉpativu cey}
  \item \emph{Recall:} say the Tamil for to enquire, then the Tamil for to register, then say \emph{muṉpativu cey} again
\end{itemize}

\begin{tcolorbox}[breakable,colback=teal!4,colframe=teal!35!black,title={Before you move on}]
What does \ta{முன்பதிவு} \ta{செய்} mean? (\textquotedblleft{}To book in advance, to reserve\textquotedblright{}.) Say it once more.
\end{tcolorbox}
\section[\texorpdfstring{\ta{விண்ணப்பி} (viṇṇappi)}{விண்ணப்பி (viṇṇappi)}]{\ta{விண்ணப்பி} (viṇṇappi) — to apply (for a job, a loan)}

\label{lesson:TA-C227-vinnappi}



\begin{quote}
Before the new one: say the Tamil for to register, then the Tamil for to book in advance.
\end{quote}

\subsection*{You'll want to know: \ta{விண்ணப்பி}}
\textbf{\ta{விண்ணப்பி}} — \emph{viṇṇappi} — \textquotedblleft{}to apply (for a job, a loan)\textquotedblright{}.

\begin{cousinweb}
\textbf{\ta{விண்ணப்பி}} (\emph{viṇṇappi}) is a verb made from \emph{viṇṇappam}, \textquotedblleft{}petition\textquotedblright{}, a well-established loan from Sanskrit \emph{vijñāpana}, \textquotedblleft{}a request, a making known\textquotedblright{}.
\end{cousinweb}

\begin{grammarlens}[title={What you've built}]
One more word for this chapter.
\end{grammarlens}

\subsection*{Your turn}
\begin{itemize}
\raggedright
  \item \emph{Say it:} \emph{viṇṇappi}
  \item \emph{Say it:} \emph{viṇṇappi}, once more
  \item \emph{Say it:} say \emph{viṇṇappi}
  \item \emph{Recall:} say the Tamil for to register, then the Tamil for to book in advance, then say \emph{viṇṇappi} again
\end{itemize}

\begin{tcolorbox}[breakable,colback=teal!4,colframe=teal!35!black,title={Before you move on}]
What does \ta{விண்ணப்பி} mean? (\textquotedblleft{}To apply (for a job, a loan)\textquotedblright{}.) Say it once more.
\end{tcolorbox}
